% =====================================================================
% §VI subsection: autonomous release confirmation (no self-drop consumption)
% 2026-09-04.  \input{} this from main.tex inside \section{Results}.
% Data: nmpc_test_results/mainline_ab2_manifest.csv (76 flights, 38 attempts)
% Verification script: src/scripts/gripper/verify_mainline_ab2.py
% =====================================================================
\subsection{Autonomous release confirmation without self-drop consumption}
\label{sec:autorelease}

The controller in Sec.~\ref{sec:method} issues the gripper-release command itself, so it
could trivially zero the payload terms in its own model the instant it commands a release.
Doing so, however, makes the closed loop depend on an internal event flag rather than on the
estimate, and that dependency is exactly what a real gripper failure invalidates. We
therefore test the stricter interface: the controller consumes \emph{only} the continuous
$m/s/J$ frame and its no-payload confidence, and never reacts to its own drop command.

\paragraph{Arms and protocol} Arm~A (baseline) retains event consumption: the release command
zeroes the payload geometry and resets the inner-loop scaling in the same frame. Arm~B
consumes the continuous estimate alone; a release is \emph{complete} only after the
no-payload confidence stays above \num{0.90} for \SI{1.0}{s}. The physical grasp procedure,
trajectory, and payload are matched within each pair. The software arms are intentionally
end-to-end architectures rather than a single-line ablation: arm~A uses its event-driven
geometry/gain path, whereas arm~B uses continuous estimates, inertia bootstrap, and setpoint
rate scheduling. Arm~A is therefore an operational reference, not a causal estimate of the
drop-notification bit alone. Three carrying conditions were run as paired A/B flights: W3 (\SI{0.15}{kg},
$e_y=\SI{0.10}{m}$, \SI{4}{m/s} figure-eight), W4 (\SI{0.30}{kg}, same trajectory), and W5
(\SI{0.15}{kg}, \SI{2}{m/s}). A pair was pre-registered as \emph{void} --- and re-flown ---
if either arm crashed \emph{before} the release command or failed the pre-registered payload-
evidence check ($p_{90}$ of estimated payload mass below threshold). The latter check did not
log the simulator joint-state acknowledgement and therefore cannot distinguish a physical
capture failure from estimator non-recognition. Void rates are reported separately in
Table~\ref{tab:void}, never absorbed into the main result.

\paragraph{Result} Among valid pairs, arm~B completed autonomous release confirmation in
\textbf{22 of 22} flights with no post-release instability observed
(Table~\ref{tab:complete}). The point estimate is \SI{100}{\percent}, but with $n=22$ the
two-sided \SI{95}{\percent} Clopper--Pearson interval is $[\SI{84.6}{\percent},
\SI{100}{\percent}]$: these data are consistent with a true failure rate up to
${\sim}\SI{15}{\percent}$ and must not be read as having established a zero failure rate.
W4 reached its pre-registered attempt cap at $n=6$ valid pairs rather than the intended 8
(its arm~A crashed before release in 9 of 16 attempts), so that cell carries meaningfully
less statistical power than W3 and W5.

\begin{table}[t]
\caption{Autonomous release completion and confirmation delay (arm~B, valid pairs only).
Delay is measured from the release command to the frame at which the no-payload confidence
has held above threshold for the full \SI{1.0}{s} hold, so it includes that hold by
construction. Arm~A completes by definition at zero delay and is not comparable.}
\label{tab:complete}
\centering
\begin{tabular}{@{}lcccc@{}}
\toprule
 & $n$ & completed & delay median [IQR] & delay max \\
\midrule
W3 (\SI{0.15}{kg}, \SI{4}{m/s}) & 8 & 8/8 & \SI{2.56}{s} [2.53, 2.85] & \SI{3.00}{s} \\
W4 (\SI{0.30}{kg}, \SI{4}{m/s}) & 6 & 6/6 & \SI{3.00}{s} [2.75, 3.25] & \SI{4.54}{s} \\
W5 (\SI{0.15}{kg}, \SI{2}{m/s}) & 8 & 8/8 & \SI{2.70}{s} [2.46, 2.95] & \SI{3.06}{s} \\
\midrule
pooled & 22 & 22/22 & \SI{2.70}{s} [2.55, 2.99] & \SI{4.54}{s} \\
\bottomrule
\end{tabular}
\\[2pt]
\footnotesize Pooled completion \SI{100}{\percent}, two-sided \SI{95}{\percent}
Clopper--Pearson CI $[\SI{84.6}{\percent}, \SI{100}{\percent}]$.
\end{table}

\paragraph{Mechanism verification} Three internal signals were logged to confirm that the
completion arises from the intended mechanism rather than from incidental decay. In all 22
valid arm~B flights the payload-presence state machine made \emph{exactly two} transitions
(one autonomous \textsc{empty}$\to$\textsc{loaded} at grasp, one \textsc{loaded}$\to$
\textsc{empty} triggered by the residual drop detector), with zero spurious transitions
during the carrying phase; the published first mass moment $s_{\text{out}}$ decayed
monotonically to zero over all 100 logged frames with no step; and the confidence crossed
threshold in every flight (\numrange{0.997}{0.999} at completion). This matters because the
\emph{internal} moment estimate $\hat s$ does \emph{not} vanish after release --- it keeps
oscillating at a level equivalent to a \SIrange{3}{16}{mm} CoM offset --- so a confidence
threshold alone could not have produced these completions; the published target must be
switched, not merely re-thresholded.

\paragraph{Post-batch detector audit} The 22/22 result validates the continuous interface and
the closed-loop recovery once the then-current residual path declares release; it is not, by
itself, a distributional validation of the detector. Five follow-on runs exposed that the
release ratio and loaded-flight ratio distributions approach within only about ten percent.
We therefore reorganized the detector around three independent sources: moment collapse,
payload quantity, and a directional thrust residual. Mass and $\Delta J$ are explicitly one
source because $\Delta J$ is algebraically derived from mass; treating them as two votes would
manufacture confidence. Absolute $\lVert s\rVert$ is only a sanity bound: its loaded minimum
(\SI{0.0039}{\kilo\gram\metre}) is lower than its post-release maximum
(\SI{0.0049}{\kilo\gram\metre}), hence it has no standalone separating power.

Table~\ref{tab:release-replay} reports the resulting historical replay of the audited detector
revision. The ratio threshold
$\rho<0.30$ was selected by leave-one-run-out cross-validation, not by the final test run, and
the same pure decision function is used online and offline. This archive produced no false
release in 2614 logged loaded frames. However, 14 of 72 releases are missed (11
of 59 after excluding manifest-invalid flights). The adversarial cross-combination of two
observed boundary values produces $\rho=1.52$, which no useful ratio threshold can accept.
Those cases require the residual source; at the time of this audit the M0 profile exposed only
the shared \SI{4.4}{N} schedule-confirmation threshold, while a \SI{0.15}{kg} release changes
weight by only \SI{1.47}{N}. The subsequent implementation separated a release-only residual
vote from that threshold. The randomized smoke below shows why the archival zero-false-release
count was insufficient to validate its fast branch.

\begin{table}[t]
\caption{Offline audit of the multi-source release decision on 73 historical SITL
runs. One run had no release command; 72 enter the release-recall denominator.}
\label{tab:release-replay}
\centering
\begin{tabular}{@{}lc@{}}
\toprule
check & result \\
\midrule
LOOCV threshold selection & 59/59 folds select $\rho=0.30$ \\
release replay & 58/72 detected; 14/72 missed \\
valid-flight subset & 48/59 detected; 11/59 missed \\
loaded-flight false releases & 0/2614 frames \\
adversarial observed-bound combination & $\rho=1.52$ (miss) \\
median detected-release delay & \SI{2.1}{s} \\
\bottomrule
\end{tabular}
\end{table}

\paragraph{Randomized fast-path smoke} The release-specific detector was scheduled for 12
repetitions at \SI{0.15}{kg}, \SI{0.10}{m} eccentricity, and a \SI{4}{m/s} figure-eight. We
stopped at 9 because fastA had already produced a flight-level false-release rate of 3/9
(Table~\ref{tab:fasta-smoke}). In a representative premature frame,
$\rho=0.694$ and the moment channel correctly remained false, but payload quantity and the
ordinary residual were simultaneously false-positive
($\hat m_p=\SI{-0.0624}{kg}$), satisfying fastA. Across the loaded segments,
$\hat m_p$ had fifth percentile \SI{-0.1032}{kg}; its longest continuous interval below the
\SI{0.03}{kg} empty threshold had median \SI{14}{s} and maximum \SI{48}{s}. Persistence and
exclusion of exact lower-bound contact therefore do not make quantity reliable during this
manoeuvre. Before the first false declaration in each flight, the moment ratio never fell below
0.789, a factor 2.6 above its 0.30 threshold.

Run~9 also exposes a stateful cascade. Its first fastA error switched the published $s$ target
to zero; because the physical payload remained attached, the presence state later re-armed,
after which the corrupted lifecycle admitted further slow/fastA declarations. These repeated
events are not independent failures: the safety statistic is 3/9 affected flights, with five
premature declarations inside them. The slow declaration at $-24.9$~s followed five consecutive
MHE solve failures: re-anchoring intentionally initialized $\hat s$ to zero, and the detector
counted this unhealthy initialization before a fresh successful solution. Thus it identifies a
missing health/freshness gate, not evidence for a physical ``near-zero'' ratio class.

\begin{table}[t]
\caption{Randomized smoke of the residual-plus-quantity fastA branch. Times are estimator
release declarations relative to the planned physical drop command.}
\label{tab:fasta-smoke}
\centering
\begin{tabular}{@{}lcc@{}}
\toprule
runs & declarations per run & timing relative to command \\
\midrule
1--5, 7 & 1 & $+0.8$ to $+1.1$ s \\
6 & 2 & $-18.1$, $+1.1$ s \\
8 & 2 & $-39.9$, $+1.0$ s \\
9 & 4 & $-51.1$, $-24.9$, $-4.8$, $+1.2$ s \\
\midrule
flight-level safety failure & \multicolumn{2}{c}{3/9} \\
\bottomrule
\end{tabular}
\end{table}

The evidence therefore supports only moment-gated decisions: slow
(quantity plus moment collapse) and fastB (strong residual plus moment collapse). We reject
fastA. The historical $\rho=1.52$ value is a synthetic cross-combination of observed boundary
values rather than evidence strong enough to justify retaining a branch with demonstrated
premature release. If a moment-gated detector times out, the safe outcome is an unresolved
state and a hover/land abort; timeout must not be promoted to empty-payload evidence. A replay
after deleting fastA reported the run-9 slow declaration again, but the source log was generated
under the old branch and is therefore intervention-contaminated rather than a clean
counterfactual. The actionable defect is still real: moment evidence must be disabled, and its
persistence reset, whenever the MHE solution is unhealthy, stale, failed, or freshly re-anchored.

\paragraph{Corrected fastA-free smoke} We removed fastA, added the estimator-health gate, and
repeated 12 online flights under the same acceptance rule: no declaration while loaded, followed
by either a valid release or a safe unresolved outcome. All 12 met that rule. Eleven produced one
post-command declaration each (nine fastB and two slow), with median detector delay
\SI{1.9}{s}; no flight declared release before the command. Three crossed the \SI{12}{s}
unresolved timeout. Two later received valid moment evidence and resolved at \SI{14.8}{s} and
\SI{14.2}{s}; one remained unresolved and, as designed, did not switch the $s$ target, clear the
model, reset adaptive state, or mark the gripper dropped.

The terminal unresolved run also separates safety from coverage. An internal residual attachment
event transitioned the presence state to \textsc{loaded}, but $\hat m_p$ never exceeded
\SI{0.09}{kg} for the 20 frames required by a legacy mass-arm latch. Because the unified detector
is still nested under that latch, it emitted no moment diagnostics at all. This is an arming-gate
coupling defect, not a release false positive: the unified detector should use the autonomous
load latch and health gate, while the legacy mass-only path retains its own mass arm. In the same
flight, mass returned close to the empty value after drop whereas raw $\hat s$ remained near
\SI{0.018}{kg.m}; forcing it to zero from the controller's drop command would violate the eventless
premise. It therefore remains an unresolved observability miss.

We corrected that arm-gate coupling and began a new acceptance batch, but stopped after four
flights when the expanded observation window exposed contamination of the running maximum
$\lVert s\rVert_{\mathrm{peak}}$. The excursions were not attachment transients: in both cohorts
they began \SIrange{26}{28}{s} after autonomous attachment, during the figure-eight ramp. Values
above the known simulation envelope of \SI{0.039}{kg.m} occurred in 49/644 logged frames (7.6\%)
of the preceding 11-flight set and 49/268 (18.3\%) of the four post-gate flights. The gate fix
therefore increased exposure to a pre-existing dynamic-estimation defect rather than causing it.
Because an inflated denominator can make an ordinary loaded moment appear collapsed, these four
flights are diagnostic, not an acceptance result. The next detector must replace the running
maximum with a robust reference formed during a healthy, low-dynamic, physically feasible window
and frozen before dynamic flight. Out-of-envelope moments are estimator-quality events, not
reference candidates.

\paragraph{Performance} Conditioned on a successful grasp and on surviving to the release
command, the paired differences are small (Table~\ref{tab:autoperf}). Steady-state tracking
RMS differs by ${\le}\SI{0.8}{mm}$ in every cell, and peak attitude, angular rate and yaw-torque
saturation are mixed in sign. We did not pre-register an equivalence test, so we claim only
that \emph{no performance degradation of engineering significance was observed}, not that the
two arms are statistically indistinguishable. The one systematic difference is the
confirmation delay of Table~\ref{tab:complete}: \SIrange{2.6}{3.0}{s} of median latency is the
price of refusing the event shortcut.

\begin{table}[t]
\caption{Tracking and attitude metrics on valid pairs, median [IQR]. All quantities are
conditioned on successful attachment and on pre-release survival; see Table~\ref{tab:void}.}
\label{tab:autoperf}
\centering
\footnotesize
\begin{tabular}{@{}llcc@{}}
\toprule
 & metric & arm~A & arm~B \\
\midrule
\multirow{3}{*}{W3 ($n{=}8$)}
 & steady RMS (\si{m})      & 0.116 [0.115, 0.118] & 0.117 [0.116, 0.119] \\
 & post-release peak (\si{m}) & 0.221 [0.206, 0.237] & 0.226 [0.221, 0.238] \\
 & peak tilt (\si{\degree})  & 12.6 [12.4, 12.7] & 12.5 [12.4, 12.5] \\
\midrule
\multirow{3}{*}{W4 ($n{=}6$)}
 & steady RMS (\si{m})      & 0.120 [0.118, 0.124] & 0.120 [0.119, 0.121] \\
 & post-release peak (\si{m}) & 0.223 [0.214, 0.238] & 0.249 [0.244, 0.255] \\
 & peak tilt (\si{\degree})  & 12.9 [12.7, 12.9] & 12.9 [12.4, 13.3] \\
\midrule
\multirow{3}{*}{W5 ($n{=}8$)}
 & steady RMS (\si{m})      & 0.044 [0.043, 0.045] & 0.044 [0.042, 0.046] \\
 & post-release peak (\si{m}) & 0.120 [0.113, 0.128] & 0.148 [0.131, 0.164] \\
 & peak tilt (\si{\degree})  & 5.37 [5.26, 5.53] & 5.25 [5.14, 5.33] \\
\bottomrule
\end{tabular}
\end{table}

\paragraph{Void flights and survivorship} Of 76 flights, 18 (\SI{24}{\percent}) were voided
and re-flown (Table~\ref{tab:void}). The two arms void for different reasons and at different
rates, so the performance comparison above is explicitly \emph{conditional} and carries a
survivorship bias that the pooled numbers cannot remove. Arm~A crashed before release more
often (11/38 vs.\ 4/38; Fisher two-sided $p=\num{0.082}$, not significant, though the
direction reproduces an earlier independent batch at 7/48 vs.\ 3/48), concentrated in W4.
Arm~B instead failed the MHE payload-evidence check in 3 flights where arm~A never did; a
companion hover batch showed the same asymmetry more strongly (0/24 vs.\ 6/24, Fisher
$p=\num{0.022}$). These historical runs saved the proximity node's attach request but not the
magnetic-joint plugin's \texttt{/gripper/state} acknowledgement. We therefore cannot identify
them as physical attachment failures: they combine capture failure and estimator
non-recognition. This unresolved attachment/observability endpoint is treated separately in
Sec.~\ref{sec:limits}; it is not included in the conditional release result.

\begin{table}[t]
\caption{Void causes by arm (76 flights, 38 attempted pairs). Void flights are excluded from
Tables~\ref{tab:complete} and~\ref{tab:autoperf} and re-flown per the pre-registered rule.}
\label{tab:void}
\centering
\begin{tabular}{@{}lccc@{}}
\toprule
 & arm~A & arm~B & total \\
\midrule
flights                       & 38 & 38 & 76 \\
crash before release          & 11 &  4 & 15 \\
payload not evidenced by MHE  &  0 &  3 &  3 \\
\midrule
valid                         & 27 & 31 & 58 \\
valid \emph{pairs}            & \multicolumn{2}{c}{22} & \\
\bottomrule
\end{tabular}
\\[2pt]
\footnotesize Crash-before-release: A 11/38 vs.\ B 4/38, Fisher two-sided $p=\num{0.082}$.
Nine of arm~A's 11 crashes are in W4, which is why that cell hit its attempt cap at 6 valid
pairs.
\end{table}

\paragraph{Takeaway} Removing the controller's consumption of its own drop event did not
produce observable post-release control degradation, at the cost of a
\SIrange{2.6}{3.0}{s} autonomous confirmation delay in the paired batch. This establishes
closed-loop feasibility of the eventless interface, conditional on reaching the release and on
the detector firing. It does not validate fastA: that later branch caused premature model
clearing in 3/9 smoke flights and is rejected. Attachment observability, pre-release divergence
under aggressive manoeuvres, dynamic contamination of the moment reference, and moment-gated
detector misses remain open limitations. The corrected 12-flight smoke establishes the intended
fail-safe semantics, but neither it nor the stopped four-flight post-gate batch is a
distribution-wide robustness claim.
